%%	build-queue:
%%	
%%	¦¦¦ very first run. No .ist files yet
%%	¦¦	index, citation/bibliography or glossary changed
%%	¦	every change apart from the above mentioned. Double run for labels and toc.
%%	
%%	¦¦¦xelatex
%%	¦¦makeglossaries
%%	¦¦makeindex
%%	¦¦bibtex
%%	¦xelatex
%%	¦xelatex
%%	

\RequirePackage[l2tabu,orthodox]{nag}
\documentclass[10pt,a4paper,titlepage,twoside,german,final]{zhawreprt}

	\setdefaultlanguage{german}
\RequirePackage{lscape}
\usepackage{pgfgantt}
\usepackage{xcolor}
	\definecolor{lgray}{RGB}{250,250,250}
	\definecolor{lgreen}{RGB}{63,127,95}
	\definecolor{lred}{RGB}{127,0,85}
	\definecolor{lblue}{RGB}{42,0,255}
\usepackage{listings}
	\lstdefinestyle{basestyle}{
		basicstyle=\small\ttfamily,
		breakatwhitespace = true,
		tabsize = 4,
		frame = double,
		numbers = left,
		numbersep = 10pt,
		numberstyle = {\tiny\emptyaccsupp},
%		firstnumber = auto,
		numberblanklines = true,
		captionpos = b,
		columns = fullflexible,
		extendedchars = true,
		float = ht,
		showtabs = false,
		showspaces=false,
		showstringspaces=false,
		breaklines=true,
%		prebreak=\Righttorque
		backgroundcolor=\color{lgray},
		keywordstyle=\color{lred}\bfseries,
		commentstyle=\color{lgreen}\ttfamily,
%		morekeywords={printstr, printhexln},
		stringstyle=\color{lblue},
		xleftmargin = \fboxsep,
		xrightmargin = -6pt,
		showstringspaces=true,
	}
	\newcommand{\setlistingCSharp}{
		\lstset{
		style = basestyle,
		language = [Sharp]C,
		%otherkeywords = {*,<,>,=,;,\{,\}},
		%deletekeywords = {...},
	}}
	\newcommand{\setlistingCpp}{
		\lstset{
		style = basestyle,
		language = C++,
		%otherkeywords = {*,<,>,=,;,\{,\}},
		%deletekeywords = {...},
	}}
	\newcommand{\setlistingJava}{
		\lstset{
		style = basestyle,
		language = Java,
		%otherkeywords = {*,<,>,=,;,\{,\}},
		%deletekeywords = {...},
	}}
	\newcommand{\setlistingLaTeX}{
		\lstset{
		style = basestyle,
		language = TeX,
		%otherkeywords = {*,<,>,=,;,\{,\}},
		%deletekeywords = {...},
	}}
	\newcommand{\setlistingMatlab}{
		\lstset{
		style = basestyle,
		language = Matlab,
		otherkeywords = {methods,enumeration,properties,classdef,Sealed,Abstract},
		%deletekeywords = {...},
	}}
\usepackage{xstring}
	\newcommand{\inlist}[2]{
		\IfSubStr{,#2,}{,\arabic{#1},}{\color{lgray!95!blue}}{\color{lgray}}
	}
\usepackage{lstlinebgrd}
	\makeatletter
	\renewcommand{\lst@linebgrd}{%
	\ifx\lst@linebgrdcolor\empty\else
		\rlap{%
			\lst@basicstyle
			\color{lgray}
			\lst@linebgrdcolor{%
				\kern-\dimexpr\lst@linebgrdsep\relax%
				\lst@linebgrdcmd{\lst@linebgrdwidth}{\lst@linebgrdheight}{\lst@linebgrddepth }%
			}%
		}%
	\fi}
	\makeatother
\usepackage[space=true]{accsupp}
	\newcommand\emptyaccsupp[1]{\BeginAccSupp{ActualText={}}#1\EndAccSupp{}}
\usepackage{rotating}
\usepackage{mathptmx}
\usepackage{amssymb}
\usepackage{textcomp}
\usepackage[squaren]{SIunits}
\usepackage{amsmath}
\usepackage{amsfonts}
\usepackage{amssymb}
\usepackage[toc,page]{appendix}
\usepackage{fontspec}
	\setmainfont{Arial} % sets the roman font
	\setsansfont{Arial} % sets the sans-sérif font
	\setmonofont{Arial} % sets the monospace font
\usepackage{microtype}
\usepackage{colortbl}
\usepackage{tabularx}
\usepackage{longtable}
\usepackage{pgf,tikz}
	\usetikzlibrary{shapes}
	\usetikzlibrary{shapes.geometric}
	\usetikzlibrary{shapes.arrows}
	\usetikzlibrary{positioning}
	\usetikzlibrary{fit}
	\usetikzlibrary{calc}
	\usetikzlibrary{patterns}
\usepackage{pgfplots}
	\pgfplotsset{compat=1.13}
\usepackage{array}
\usepackage{natbib}
	\bibliographystyle{agsm}
\usepackage{usecases}
\usepackage{footnote}
	\makesavenoteenv{description}
\usepackage{makeidx}
	\makeindex
\usepackage[nottoc]{tocbibind}
\makeatletter
	\if@inltxdoc\else
	  \renewenvironment{theindex}%
	    {\if@twocolumn
	       \@restonecolfalse
	     \else
	       \@restonecoltrue
	     \fi
	     \if@bibchapter
	        \if@donumindex
	          \refstepcounter{section}
	          \twocolumn[\vspace*{2\topskip}%
	                     \@makechapterhead{\indexname}]%
	          \addcontentsline{toc}{section}{\protect\numberline{\thesection}\indexname}
	          \sectionmark{\indexname}
	        \else
	          \if@dotocind
	            \twocolumn[\vspace*{2\topskip}%
	                       \@makeschapterhead{\indexname}]%
	            \prw@mkboth{\indexname}
	            \addcontentsline{toc}{section}{\protect\numberline{\thesection}\indexname}
	          \else
	            \twocolumn[\vspace*{2\topskip}%
	                       \@makeschapterhead{\indexname}]%
	            \prw@mkboth{\indexname}
	          \fi
	        \fi
	     \else
	        \if@donumindex
	          \twocolumn[\vspace*{-1.5\topskip}%
	                     \@nameuse{\@tocextra}{\indexname}]%
	          \csname \@tocextra mark\endcsname{\indexname}
	        \else
	          \if@dotocind
	            \twocolumn[\vspace*{-1.5\topskip}%
	                       \toc@headstar{\@tocextra}{\indexname}]%
	            \prw@mkboth{\indexname}
	            \addcontentsline{toc}{\@tocextra}{\indexname}
	          \else
	            \twocolumn[\vspace*{-1.5\topskip}%
	                       \toc@headstar{\@tocextra}{\indexname}]%
	            \prw@mkboth{\indexname}
	          \fi
	        \fi
	     \fi
	   \thispagestyle{plain}\parindent\z@
	   \parskip\z@ \@plus .3\p@\relax
	   \let\item\@idxitem}
	   {\if@restonecol\onecolumn\else\clearpage\fi}
	\fi
\makeatother
\usepackage{etoolbox}
\makeatletter
	\renewcommand\listoftables{%
	    \section{Tabellenverzeichnis}%
	    \@mkboth{\MakeUppercase\listtablename}%
	        {\MakeUppercase\listtablename}%
	    \@starttoc{lot}%
	}
	\renewcommand\listoffigures{%
	    \section{Abbildungsverzeichnis}%
	    \@mkboth{\MakeUppercase\listfigurename}%
	        {\MakeUppercase\listfigurename}%
	    \@starttoc{lof}%
	}
	\renewcommand\lstlistoflistings{%
	    \section{Listingverzeichnis}%
	    \@mkboth{\MakeUppercase\listfigurename}%
	        {\MakeUppercase\listfigurename}%
	    \@starttoc{lol}%
	}
	\renewenvironment{thebibliography}[1]
	     {\section{\bibname}% <-- this line was changed from \chapter* to \section*
	      \@mkboth{\MakeUppercase\bibname}{\MakeUppercase\bibname}%
	      \list{\@biblabel{\@arabic\c@enumiv}}%
	           {\settowidth\labelwidth{\@biblabel{#1}}%
	            \leftmargin\labelwidth
	            \advance\leftmargin\labelsep
	            \@openbib@code
	            \usecounter{enumiv}%
	            \let\p@enumiv\@empty
	            \renewcommand\theenumiv{\@arabic\c@enumiv}}%
	      \sloppy
	      \clubpenalty4000
	      \@clubpenalty \clubpenalty
	      \widowpenalty4000%
	      \sfcode`\.\@m}
	     {\def\@noitemerr
	       {\@latex@warning{Empty `thebibliography' environment}}%
	      \endlist}
\makeatother
	
	\renewcommand\appendixname{Appendix}
\makeatletter
	\renewenvironment{theindex}{
		\renameindex
		\let\ps@plainorig\ps@plain
		\let\ps@plain\ps@scrheadings
		\if@twocolumn
			\@restonecolfalse
		\else
			\@restonecoltrue
		\fi
		\columnseprule \z@
		\columnsep 35\p@
		\twocolumn[\section{\indexname}]%
		\@mkboth{\MakeUppercase\indexname}{\MakeUppercase\indexname}%
		\thispagestyle{plain}\parindent\z@
		\parskip\z@ \@plus .3\p@\relax
	\let\item\@idxitem}
	{\if@restonecol\onecolumn\else\clearpage\fi}
\makeatother

\usepackage{varioref}
\usepackage[colorlinks=true, linkcolor=black, citecolor=black, plainpages=false, unicode, pdfencoding=auto ,backref=page]{hyperref}
\usepackage{cleveref}
\usepackage[automake,acronym,numberedsection]{glossaries-extra}
	\renewcommand*{\glspostdescription}{}
	\newglossary[slg]{symbols}{sym}{sbl}{Symbolverzeichnis}
	\makeglossaries
	\loadglsentries{glossaryentries.tex}
	\pdfstringdefDisableCommands{\let\textenglish\@firstofone\let\textgerman\@firstofone}
	\makeatletter
	\renewcommand*{\@@glossarysec}{section}
	\makeatother
\if false
%\newcommand{\FullGlossaryEntry}[7]{%Abkz., Typ, Name, Name(plural), Beschr., Beschr.(plural), Erkl.
%\newglossaryentry{#1}{type=\acronymtype,
%name={\textit{#3}},
%plural={\textit{#4}},
%description={\textit{#5}},
%first={\textit{#5}(nachfolgend #3)\glsadd{_#1}},
%firstplural={\textit{#6}(nachfolgend \glsentryplural{#4})},
%see=[Glossar:]{_#1}
%}
%\longnewglossaryentry{_#1}{type=#2,
%name={\textit{#5}},
%plural={\textit{#6}}
%}{\hspace{0pt}\\#7}
%}


% Abkuerzungen
\newacronym{jlu}{JLU}{Justus-Liebig-Universität}
\newacronym{hrz}{HRZ}{Hochschulrechenzentrum}
\newacronym[plural=LEDs, longplural={light-emitting diodes}]{led}{LED}{light-emitting diode}
\newacronym[plural=EEPROMs, longplural={electrically erasable programmable read-only memories}]{eeprom}{EEPROM}{electrically erasable programmable read-only memory}

%  Glossareintraege
\newglossaryentry{culdesac}{name=cul-de-sac, description={passage or street closed at one end}, plural=culs-de-sac}
\newglossaryentry{elite}{name={é}lite, description={select group or class}, sort=elite}
\newglossaryentry{elitism}{name={é}litism, description={advocacy of dominance by an \gls{elite}}, sort=elitism}
\newglossaryentry{attache}{name=attaché, description={person with special diplomatic responsibilities}}

% Eintraege für Symbolliste
\newglossaryentry{ohm}{type=symbols, name={\ensuremath{\Omega}}, sort=Ohm, symbol={\ensuremath{\Omega}}, description={unit of electrical resistance}}
\newglossaryentry{angstrom}{type=symbols, name={\AA}, sort=angström, symbol={\AA}, description={non-SI unit of length}}
%
%\FullGlossaryEntry{zhaw}{main}%
%{ZHaW}{}%
%{Zürcher Hochschule der angewandten Wissenschaften}{}%
%{Name der Hochschule meines Vertrauens.}
%
%\FullGlossaryEntry{mnmt1}{main}%
%{MNMT1}{}%
%{Mathematik: Numerik für Maschinentechnik 1}{}%
%{Erstes Numerik-Modul im Studiengang Maschinentechnik an der \gls{zhaw}. Beinhaltet Taylor- und Fourier-Reihen, Numerik gewöhnlicher Differentialgleichungen anhand des Euler-, Taylor- und Runge-Kutta-Verfahren, Numerik nichtlinearer Differentialgleichungen, Lagrange- und Newton-Interpolation, Splines und Ausgleichsrechnung.}
%
%\FullGlossaryEntry{mnmt2}{main}%
%{MNMT2}{}%
%{Mathematik: Numerik für Maschinentechnik 2}{}%
%{Zweites Numerik-Modul im Studiengang Maschinentechnik an der \gls{zhaw}. Beinhaltet Randwertprobleme, numerische Differentiation und Integration, Numerik partieller Differentialgleichungen anhand der \gls{fdm} und \gls{fem}, Stabilität von numerischen Verfahren zur Lösung von gewöhnlichen Differentialgleichungen, Schrittweitensteuerung und implizite Runge-Kutta-Verfahren.}
%
%\FullGlossaryEntry{fdm}{main}%
%{FDM}{}%
%{Finite Differenzen Methode}{}%
%{Numerisches Verfahren zur Lösung von \glspl{rwp}}
%
%\FullGlossaryEntry{fem}{main}%
%{FEM}{}%
%{Finite Elemente Methode}{}%
%{Numerisches Verfahren zur Lösung von \glspl{rwp}}
%
%\FullGlossaryEntry{rwp}{main}%
%{RWP}{RWPs}%
%{Randwertproblem}{Randwertprobleme}%
%{Ein Differentialproblem, bei dem diskrete Funktionswerte einer bestimmten Ordnung an beliebigen Stellen gegeben sind.}
%
%\FullGlossaryEntry{awp}{main}%
%{AWP}{AWPs}%
%{Anfangswertproblem}{Anfangswertprobleme}%
%{Ein Differentialproblem, bei dem alle diskrete Funktionswerte an einer Anfangsstelle gegeben sind.}
%
%\FullGlossaryEntry{dgl}{main}%
%{DGL}{DGLn}%
%{Differentialgleichung}{Differentialgleichungen}%
%{Eine Gleichung, bei der Ableitungen beliebiger Ordnung einer gesuchten Funktion vorkommen.}
%
%\FullGlossaryEntry{ode}{main}%
%{ode}{ode's}%
%{ordinary differential equation}{ordinary differential equations}%
%{Englische Bezeichnung für \gls{dgl}}
%
%
%\FullGlossaryEntry{dgls}{main}%
%{DGLS}{DGLSe}%
%{Differentialgleichungssystem}{Differentialgleichungssysteme}%
%{Ein System von Differentialgleichungen}
%
%\FullGlossaryEntry{ods}{main}%
%{ods}{ods's}%
%{ordinary differential system}{ordinary differential systems}%
%{Englische Bezeichnung für \gls{dgls}}
\fi

\logofilename{images/logos/SoE/de/de-soe-cmyk.png}
\projecttype{PA}
\major{HS16 Studiengang Informatik}
\title{[Trip]od}
\shorttitle{[Trip]od}
\author{Gruppe 15}
\authors{Gruppe 15: Fabio Costi, Raphael Emberger,\\Nicolas Kaiser, Julian Visser}
\mainreferee{Daniel Liebhart}
\referee{Liby Kunthrayil}
\industrypartner{}
\extreferee{}
\setdate{\today}

\newcommand{\AddRequirement}[2]{
\textbf{/#1#2/}
}
\newcommand{\F}[1]{
\AddRequirement{F1.}{#1}
}
\newcommand{\W}[1]{
\AddRequirement{F2.}{#1}
}
\newcommand{\D}[1]{
\AddRequirement{F3.}{#1}
}
\newcommand{\R}[1]{
\AddRequirement{F4.}{#1}
}
\newcommand{\TF}[1]{
\AddRequirement{T.}{#1}
}
\newcommand{\tableheader}[2]{\multicolumn{#1}{c}{\raisebox{-0.3em}[0ex][0ex]{\large{\textbf{#2}}}}}

\numberwithin{table}{chapter}
%\renewcommand{\thetable}{\arabic{chapter}.\arabic{table}}

\begin{document}

\makezhawtitle{
\begin{center}
\begin{tikzpicture}
\fill[fill=green!20,draw=green!80!black](0,0) rectangle (12,5);
\node[text=white,rotate=30] at (6,2.5){\Huge\textbf{[Trip]od}};
\end{tikzpicture}
\end{center}
}

\tableofcontents

\chapter{Einleitung}\label{chp:Introduction}
Unsere Gruppe hat beim Brainstorming festgestellt, dass es noch kein geeignetes Werkzeug gibt, mit welchem man seine Ferien planen kann, sämtliche Informationen an einem Ort gesammelt sind und auf diese Informationen von jedem Mitglied der Reisegruppe jederzeit zugegriffen werden kann. Aus diesem Grund haben wir beschlossen, diese Lücke zu schliessen.\\
Geplant ist in unserem Projekt eine Webapplikation, mit welcher der Nutzer bequem seine künftigen Ferien planen kann.\\
Der Ferienplaner trägt den Namen \gls{tripod}, welcher sich aus folgenden drei Punkten zusammensetzt:
\begin{itemize}
\item TRIP: steht für die kommende Reise
\item POD: englisches Sprichwort "like two peas in a pod" (übersetzt: wie zwei Erbsen in einer Schote), steht in unserem Fall dafür, dass die Reisenden alle eine gemeinsame Faszination teilen
\item TRIPOD: das Stativ, auf welchem das Fernrohr befestigt ist, um auf der Reise in die Ferne zu schauen
\end{itemize}
Der Nutzer kann sich auf der Webseite registrieren und sich einloggen. Anschliessend kann die Planung gestartet werden. Man legt zuerst die Ferien an und kann, wenn dies gewünscht ist, beliebig Freunde dazu einladen. Diese Freunde müssen sich auch registrieren und können dann bei der Planung mitwirken. Im nächsten Schritt einigt man sich auf einen Zeitraum, in dem die Ferien stattfinden, auf ein Reiseziel und ein Budget.\\
Nun informiert sich jeder Teilnehmer über mögliche Verkehrsmittel, um zu der gewünschten Destination zu gelangen und über mögliche Unterkünfte. Diese Suche findet ausserhalb unserer Webseite statt. Ein Wunschziel von uns ist es, dass wir über verschiedene \gls{api}'s der Anbieter (Schnittstellen zu deren Webseiten) die Suche direkt auf unserer Webseite ermöglichen können.\\
Wenn ein Nutzer ein attraktives Angebot für ein Verkehrsmittel oder eine Unterkunft gefunden hat, kann er dieses auf unserer Seite eintragen. Dieses Angebot ist dann für jeden Teilnehmer einsehbar. Diese Angebote können dann in der Gruppe mittels der Kommentar-Funktion besprochen werden. Für das definitive Buchen der Reise bietet unsere Webseite keine Funktion an. Dies geschieht weiterhin über die jeweiligen Anbieter.\\
Da während der Reise oft kein Internet verfügbar ist, bieten wir eine Funktion an, mit welcher sich alle Informationen exportieren lassen und sie somit jederzeit auch offline zur Verfügung stehen.

\chapter{Zielbestimmungen}\label{chp:DefinitionOfGoals}
\section{Musskriterien}\label{sec:MustCriteria}

In den folgenden Tabellen sind die Musskriterien des \gls{tripod}'s definiert.
\begin{table}[ht]\centering
\begin{tabular}{l}\hline
\tableheader{1}{Account verwalten}\\[0.3em]\hline
\textbf{Kriterium}\\\hline
Account anlegen\\\hline
Account bearbeiten (Passwort ändern, persönliche Daten ausfüllen)\\\hline
Benutzer an- und abmelden\\\hline
Passwort vergessen\\\hline
\end{tabular}
\caption{Musskriterien Account}\label{tbl:MustAccount}
\end{table}

\begin{table}[ht]\centering
\begin{tabular}{l}\hline
\tableheader{1}{Ferien verwalten}\\[0.3em]\hline
\textbf{Kriterium}\\\hline
Ferien anlegen\\\hline
Personen hinzufügen und entfernen\\\hline
Reisedaten festlegen\\\hline
Budget festlegen\\\hline
Destination hinzufügen, bearbeiten, entfernen \\\hline
Verkehrsmittel hinzufügen, bearbeiten, entfernen\\\hline
Unterkunft hinzufügen, bearbeiten, entfernen\\\hline
Kommentar-Funktion\\\hline
Offline-Export\\\hline
\end{tabular}
\caption{Musskriterien Ferien}\label{tbl:MustVacation}
\end{table}\newpage
\section{Wunschkriterien}\label{sec:WishCriteria}
Um den Funktionsumfang bei verbleibender Zeit zu erweitern, wurden die in der folgenden Tabelle stehenden Wunschkriterien definiert.
\begin{table}[ht]\centering
\begin{tabular}{l}\hline
\tableheader{1}{Ferien verwalten}\\[0.3em]\hline
\textbf{Kriterium}\\\hline
Dokumente anhängen\\\hline
Mehrere Destinationen\\\hline
API-basierte Flugsuche\\\hline
API-basierte Unterkunftssuche\\\hline
\end{tabular}
\caption{Wunschkriterien Ferien}\label{tbl:WishVacation}
\end{table}
\section{Abgrenzungskriterien}\label{sec:DistinctionCriteria}
Die in diesem Kapitel definierten Abgrenzungskriterien sollen bewusst nicht erreicht werden.
\begin{table}[ht]\centering
\begin{tabular}{l|p{2.5cm}|p{5cm}|p{2cm}|p{1.5cm}}\hline
\tableheader{5}{Allgemein}\\[0.3em]\hline
\textbf{Kriterium} & \textbf{Kurzbeschrieb} & \textbf{Beschreibung} & \textbf{Komplexität} & \textbf{Priorität}\\\hline
\D{1000} & Booking & Die tatsächliche Tätigung von Käufen oder Buchungen soll nicht unterstützt werden. & hoch & tief\\\hline
\end{tabular}
\caption{Abgrenzungskriterien Allgemein}\label{tbl:DistinctionGeneral}
\end{table}
\chapter{Produkteinsatz}\label{chp:ProductApplication}
\section{Anwendungsbereiche}\label{sec:FieldOfApplience}
Genutzt wird der Ferienplaner über das Web. Dies ermöglicht, dass mehrere Personen problemlos zusammen ihre Reise planen können.\\
Theoretisch kann jeder Mensch, der Zugang zum Internet hat, unsere Applikation benutzen, jedoch ist vorerst nur eine deutsche Webseite geplant, sodass in einem ersten Schritt auf deutschsprachige Nutzer gezielt wird.
\section{Benutzergruppen des Produktes}\label{sec:TargetAudience}
Das Produkt gilt der öffentlichen, nicht-kommerziellen Nutzung und ist für Gruppen wie auch Einzelpersonen geeignet. Dies umfasst Personen jeden Geschlechtes oder Alters. Ein wichtiger Punkt unserer Webseite ist auch das Budget, da gerade wir Studenten oftmals nicht über grosse, finanzielle Möglichkeiten verfügen. Somit ist unsere Applikation bestens geeignet für reiselustige Personen, welche gerne ihre Reise selber planen oder sich den Luxus eines Reisebüros nicht leisten können.

\chapter{Funktionale Anforderungen}\label{chp:FunctionalRequirements}
Der Ferienplaner ist über einen eigenen Account zu verwalten. Jeder Reisende braucht einen eigenen Account (Dummy-Accounts für Gruppen könnten eine Alternative sein). Mit diesem Account lassen sich neue Ferien anlegen, denen andere Accounts hinzugefügt werden können. Bei Gruppen ist der Gruppenersteller automatisch der Gruppenadmin. Ferien können nun terminlich eingeschränkt, Abreise- und Rückkehrdaten festgelegt und Reiseziele bestimmt werden. Weitere Einstellmöglichkeiten (Minimum-Budget, Maximum-Komfort, etc.) können ebenfalls bestimmt werden. All diese Parameter können vom Admin uneingeschränkt und von Gruppenmitgliedern nach Revision des Admins und anderer Gruppenmitglieder aufgenommen werden.\\
Nun kann auf Knopfdruck die Berechnung gestartet werden. Dabei werden über die API's von mehreren Anbietern Angebote gesammelt und Rechnungen getätigt. Der Output soll aus den Kosten der Ideallösung (Auswahl der Möglichkeiten evtl. ebenfalls möglich), einer Tabelle der buchbaren Mittel und einem Zeitplan bestehen. Das effektive Buchen geht dann aus Sicherheitsgründen weiterhin händisch über den Admin.

\begin{table}[ht]\centering
\begin{tabular}{l|p{2.5cm}|p{5cm}|p{2cm}|p{1.5cm}}\hline
\tableheader{5}{Account verwalten}\\[0.3em]\hline
\textbf{Kriterium} & \textbf{Kurzbeschrieb} & \textbf{Beschreibung} & \textbf{Komplexität} & \textbf{Priorität}\\\hline
\F{1000} & Account anlegen & Es kann ein Account mit den folgenden Daten angelegt werden:\linebreak
\begin{itemize}
\item E-Mail Adresse [zwingend]
\item Passwort [zwingend]
\item Name (Vor- und Nachname) [zwingend]
\item Telefonnummer [fakultativ]
\end{itemize}
& mittel & hoch\\\hline
\F{1010} & Account bearbeiten & Das Passwort und die persönlichen Daten können angepasst werden. Einzig die E-Mail Adresse kann nicht verändert werden. & tief & hoch\\\hline
\F{1020} & Benutzer anmelden & Ein Benutzer kann sich mit seinem Account (E-Mail und Passwort) anmelden. & mittel & hoch\\\hline
\F{1030} & Benutzer abmelden & Ein Benutzer kann sich, wenn er angemeldet ist, auch wieder abmelden. & tief & hoch\\\hline
\F{1040} & Passwort vergessen & Wenn ein Benutzer sein Passwort vergessen hat, kann er mittels Eingabe seiner E-Mail Adresse ein neues Passwort anfordern. & tief & hoch\\\hline
\end{tabular}
\caption{Funktionale Anforderungen Account}\label{tbl:FuncAccount}
\end{table}

\begin{table}[ht]\centering
\begin{tabular}{l|p{2.5cm}|p{5cm}|p{2cm}|p{1.5cm}}\hline
\tableheader{5}{Ferien verwalten}\\[0.3em]\hline
\textbf{Kriterium} & \textbf{Kurzbeschrieb} & \textbf{Beschreibung} & \textbf{Komplexität} & \textbf{Priorität}\\\hline
\F{2000} & Ferien anlegen & Jeder Benutzer kann Ferien erstellen. Eine Ferienplanung beinhaltet folgende Informationen:\linebreak
\begin{itemize}
\item Titel [zwingend]
\item Beschreibung [zwingend]
\item Ersteller [zwingend]
\item Personen [fakultativ, kann auch nur der "Ersteller" sein]
\item Reisedaten (An- und Abreisedatum) [zwingend]
\item Budget [fakultativ]
\item Destinationen [mindestens Eine]
\item Verkehrsmittel [fakultativ]
\item Unterkünfte [fakultativ]
\item Kommentare [fakultativ]
\end{itemize}
& tief & hoch\\\hline
\F{2010} & Personen hinzufügen & Jeder Benutzer kann zusätzliche Personen (registrierte Benutzer) per E-Mail Adresse zur Ferienplanung hinzufügen. Die hinzugefügte Person wird mittels E-Mail über die neu verfügbare Ferienplanug informiert. & tief & hoch\\\hline
\F{2020} & Personen entfernen & Jeder Benutzer kann sich oder andere Personen (registrierte Benutzer) entfernen. Die entfernte Person wird mittels E-Mail über den Ausschluss aus der Ferienplanung informiert. & tief & hoch\\\hline
\F{2030} & Reisedaten festlegen & Der Zeitraum der Ferien kann festgelegt und verändert werden. & tief & hoch\\\hline
\F{2040} & Budget festlegen & Das Budget kann festgelegt und verändert werden. Dieses besteht aus:
\begin{itemize}
\item Budget für die Reise
\item Budget für die Unterkunft
\end{itemize}
& tief & hoch\\\hline
\F{2050} & Destination hinzufügen & Jeder Benutzer kann eine Destination hinzufügen. & tief & hoch\\\hline
\F{2060} & Destination bearbeiten & Jeder Benutzer kann eine Destination bearbeiten. & tief & hoch\\\hline
\F{2070} & Destination entfernen & Jeder Benutzer kann eine Destination entfernen. & tief & hoch\\\hline
\F{2080} & Verkehrsmittel hinzufügen & Jeder Benutzer kann ein Verkehrsmittel hinzufügen. Hierzu kann ein Link des jeweiligen Angebotes hinterlegt werden. & tief & hoch\\\hline
\F{2090} & Verkehrsmittel bearbeiten & Jeder Benutzer kann ein Verkehrsmittel bearbeiten. & tief & hoch\\\hline
\F{2100} & Verkehrsmittel entfernen & Jeder Benutzer kann ein Verkehrsmittel entfernen. & tief & hoch\\\hline
\F{2110} & Unterkunft hinzufügen & Jeder Benutzer kann eine Unterkunft hinzufügen. Hierzu kann ein Link des jeweiligen Angebotes hinterlegt werden. & tief & hoch\\\hline
\F{2120} & Unterkunft bearbeiten & Jeder Benutzer kann eine Unterkunft bearbeiten. & tief & hoch\\\hline
\F{2130} & Unterkunft entfernen & Jeder Benutzer kann eine Unterkunft entfernen. & tief & hoch\\\hline
\F{2140} & Kommentar-Funktion & Jeder Benutzer kann Kommentare hinterlassen (Chat). & mittel & hoch\\\hline
\F{2150} & Offline-Export & Jeder Benutzer kann einen Export mit den wichtigsten Daten der Ferien herunterladen. & mittel & hoch\\\hline
\W{2000} & Dokumente anhängen & Booking-Bestätigungen und ähnliche Dokumente sollen abgelegt werden können. & hoch & tief\\\hline
\W{2010} & Mehrere Destinationen & Es soll möglich sein, eine Planung für mehrere Destinationen hintereinander zu machen. & hoch & tief\\\hline
\W{2020} & API-basierte Flugsuche & Es soll möglich sein, Flugangebote über eine API abzurufen. z.B. Google Flights & hoch & tief\\\hline
\W{2030} & API-basierte Unterkunftssuche & Es soll möglich sein, Unterkünfte über eine API abzurufen. z.B. Booking.com & hoch & tief\\\hline
\end{tabular}
\caption{Funktionale Anforderungen Ferien}\label{tbl:FuncVacation}
\end{table}\newpage

\chapter{Funktionsbaum}\label{chp:FunctionTree}
\begin{minipage}{\textwidth}
%\begin{landscape}
\definecolor{MainBack1}{RGB}{66,66,68}
\definecolor{MainBack2}{RGB}{86,86,88}
\definecolor{MainFront}{RGB}{220,220,220}
\definecolor{SubBack1}{RGB}{240,240,240}
\definecolor{SubBack2}{RGB}{255,255,255}
\definecolor{SubFront}{RGB}{46,46,48}
\newcommand{\NewMainNode}[3]{ %NewNodeName, Pos, Text
\node(#1)[MainNode] at (#2) {#3};
}
\newcommand{\AddMainNode}[5]{ %NewNodeName, Parent, XOff, YOff, Text
\node(#1)[MainNode,below=of #2,shift={(#3, #4)}] {#5};
\coordinate(Temp) at ($(#1)-(#3, -1.3cm)$);
\draw[line](#2) -- (Temp) -| (#1);
}
\newcommand{\CreateNodeOffset}[2]{ %NewNodeName, Parent
\path (#2) + (30pt, -20.5pt) coordinate (#1);
}
\newcommand{\NewSubNode}[3]{ %NewNodeName, Parent, Text
\node(#1)[SubNode,below=of #2,shift={(40pt, 0pt)}] {#3};
\draw[line](#2) ++ (SubOffset) |- (#1);
}
\newcommand{\AddSubNode}[4]{ %NewNodeName, Parent, BelowSubNode Text
\node(#1)[SubNode,below=of #3] {#4};
\draw[line](#2) ++ (SubOffset) |- (#1);
}
\begin{turn}{90}
\begin{tikzpicture}[
text height = 0.1cm,
every node/.style = {
	column sep = 0.5cm,
	row sep = 1cm,
	thick,
	node distance = 1em,
	align = center,
	font = \small
},
base/.style = {
	text centered,
	inner sep = 1pt,
	minimum height = 30pt,
	minimum width = 100pt,
},
MainNode/.style = {
	base,
	rectangle,
	bottom color = MainBack1,
	top color = MainBack2,
	text = MainFront,
	draw = MainBack2
},
SubNode/.style = {
	base,
	rectangle,
	bottom color = SubBack1,
	top color = SubBack2,
	text = SubFront,
	draw = SubFront
},
line/.style = {
	draw,
	-stealth,
	shorten > = 2pt,
	thick
}
]

\coordinate(SubOffset) at (-30pt,-15pt);
\coordinate(MainOffset) at (12cm, 0cm);

%% LAYOUT
\NewMainNode{F2}{0, 0}{Ferien verwalten}

\NewSubNode{F2000}{F2}{Ferien anlegen}
\AddSubNode{F20100}{F2}{F2000}{Personen verwalten}

\node(F2010)[SubNode,right=of F20100,shift={(1cm, 0pt)}] {Personen hinzufügen};
\draw[line](F20100) |- (F2010);

\node(F2020)[SubNode,below=of F2010] {Personen abmelden};
\draw[line](F20100) ++ (2.5cm, 0pt) |- (F2020);

\AddSubNode{F2020}{F2}{F20100}{Daten festlegen}
\AddSubNode{F2030}{F2}{F2020}{Budget festlegen}
\AddSubNode{F20400}{F2}{F2030}{Destinationen verwalten}

\node(F2040)[SubNode,right=of F20400,shift={(11cm, 0pt)}] {Destinationen hinzufügen};
\draw[line](F20400) |- (F2040);
\node(F2050)[SubNode,below=of F2040] {Destinationen bearbeiten};
\draw[line](F20400) ++ (12.5cm, 0pt) |- (F2050);
\node(F2060)[SubNode,below=of F2050] {Destinationen entfernen};
\draw[line](F20400) ++ (12.5cm, 0pt) |- (F2060);

\AddSubNode{F20700}{F2}{F20400}{Verkehrsmittel verwalten}

\node(F2070)[SubNode,right=of F20700,shift={(6cm, 0pt)}] {Verkehrsmittel hinzufügen};
\draw[line](F20700) |- (F2070);
\node(F2080)[SubNode,below=of F2070] {Verkehrsmittel bearbeiten};
\draw[line](F20700) ++ (7.5cm, 0pt) |- (F2080);
\node(F2090)[SubNode,below=of F2080] {Verkehrsmittel entfernen};
\draw[line](F20700) ++ (7.5cm, 0pt) |- (F2090);

\AddSubNode{F21000}{F2}{F20700}{Unterkunft verwalten}

\node(F2100)[SubNode,right=of F21000,shift={(1cm, 0pt)}] {Unterkunft hinzufügen};
\draw[line](F21000) |- (F2100);
\node(F2110)[SubNode,below=of F2100] {Unterkunft bearbeiten};
\draw[line](F21000) ++ (2.5cm, 0pt) |- (F2110);
\node(F2120)[SubNode,below=of F2110] {Unterkunft entfernen};
\draw[line](F21000) ++ (2.5cm, 0pt) |- (F2120);

\AddSubNode{F2130}{F2}{F21000}{Kommentar-Funktion}
\AddSubNode{F2140}{F2}{F2130}{Dokumente anhängen}


\NewMainNode{F1}{$(F2)+(MainOffset)$}{Account}
\NewSubNode{F1010}{F1}{Account anlegen}
\AddSubNode{F1020}{F1}{F1010}{Account bearbeiten}

\end{tikzpicture}
\end{turn}
%\end{landscape}
\end{minipage}
\chapter{Nicht-Funktionale Anforderungen}\label{chp:NonFunctionalRequirements}
\begin{table}[ht]\centering
\begin{tabular}{l|p{3.0cm}|p{5cm}|p{2cm}|p{1.25cm}}\hline
\textbf{Kriterium} & \textbf{Kurzbeschrieb} & \textbf{Beschreibung} &\textbf{Komplexität} & \textbf{Priorität}\\\hline
\R{1000} & Gebrauchsfähigkeit (Usability) & Die Webseiten müssen durch den Benutzer, welcher dem Benutzerprofil entspricht, ohne weitere Hilfe verwendet werden können. & mittel & hoch\\
\R{1010} & Fehlertoleranz & Hinweise und Fehlermeldungen müssen für den Benutzer verständlich formuliert sein und eine Hilfestellung zur Problemlösung beinhalten. & mittel & mittel\\
\R{1020} & Sprache \& länderspezifische Einstellungen & Die Webseiten sind in deutscher Sprache (Schweiz) verfasst, verwenden den Zeichensatz UTF-8 und die Schweiz-spezifischen Einstellungen von Datum, Zeit, Zahlen und Währung. & tief & hoch\\
\R{1030} & Zielplattform (Server) & Die Web-Applikation muss als JavaServer-Pages auf dem zur Verfügung gestellten virtuellen Server unter Verwendung einer SQL-Datenbank mit Apache Tomcat betrieben werden. & mittel & hoch\\
\R{1040} & Zielplattform (Client) & Die Webseiten werden in der aktuellsten freigegebenen Version des Mozilla Firefox und Google Chrome korrekt dargestellt. & mittel & hoch\\
\R{1050} & Werkzeuge zur Entwicklung & Als Projektmanagement-Tool und zur Verwaltung des Sourcecode und der Dokumente muss der zur Verfügung gestellte github Server verwendet werden. & tief & hoch\\
\R{1060} & Robustheit & Auch nach einem Neustart des virtuellen Servers muss die Webseite voll funktionsfähig sein. & mittel & hoch\\
\R{1070} & Testbarkeit & Für die Durchführung der Tests und der Abnahme müssen sinnvolle Testdaten in genügendem Umfang zur Verfügung gestellt werden. & hoch & hoch\\
\end{tabular}
\caption{Nicht-Funktionale Anforderungen}\label{tbl:NonFunctionalRequirements}
\end{table}
\chapter{Abnahmekriterien Anforderungen}\label{chp:TestRequirements}
\begin{table}[ht]\centering
\begin{tabular}{l|p{4cm}|p{8cm}}\hline
Kriterium & Kurzbeschrieb & Beschreibung\\\hline
\TF{1000} & Account anlegen und Profil befüllen & Es ist ein Account anzulegen mit dem Namen "Test Account" und der E-Mail-Addresse "account@test.me". Weiter soll das Profil wie folgt befüllt werden:\ldots
\end{tabular}
\caption{Abnahmekriterien}\label{tbl:TestRequirements}
\end{table}
\chapter{Verzeichnisse}\label{chp:Index}
\printglossaries\label{sec:Glossar}
\listoftables
\printindex
\newpage
\bibliography{reference}\label{sec:Bibliography}
\end{document}